\documentclass[11pt,a4paper]{article}

\usepackage[T1]{fontenc}
\usepackage[utf8]{inputenc}
\usepackage{lmodern}
\usepackage{geometry}
\usepackage{amsmath,amssymb,amsthm}
\usepackage{bm}
\usepackage{booktabs}
\usepackage{hyperref}
\usepackage{enumitem}

\geometry{margin=2.5cm}
\setlist[itemize]{itemsep=0.25em, topsep=0.4em}
\setlist[enumerate]{itemsep=0.25em, topsep=0.4em}

\title{Mathematical Notes For The HydroModPy Boussinesq Backend}
\author{HydroModPy}
\date{\today}

\begin{document}

\maketitle
\tableofcontents

\section{Purpose Of This Note}

This note gives a mathematical reading of the Boussinesq backend implemented in
\texttt{hydromodpy/solver/boussinesq}. It is not a full research-paper style
derivation. Its role is more practical:

\begin{itemize}
\item define the unknowns and the sign conventions used by the code,
\item show how the finite-volume residual is assembled,
\item explain the additional regularized source terms currently present in the
      implementation,
\item connect the equations to the main Python modules.
\end{itemize}

The intended audience is a developer or modeller who wants to understand what
the backend solves and how that solve is encoded.

\section{Geometric Setting And Unknowns}

The solver works on a 2D triangular mesh of aquifer cells. Each cell $i$ has:

\begin{itemize}
\item area $A_i$,
\item top elevation $z^{\text{top}}_i$,
\item bottom elevation $z^{\text{bot}}_i$,
\item hydraulic conductivity $K_i$,
\item storage coefficient $S_i$.
\end{itemize}

The primary unknown is the cell-centered hydraulic head
\[
h_i.
\]

From $h_i$, the code reconstructs the saturated thickness
\[
b_i(h) = \min\!\left(
    \max\!\left(h_i - z^{\text{bot}}_i, 0\right),
    z^{\text{top}}_i - z^{\text{bot}}_i
\right).
\]

This clipping is important. It means the nonlinear iterate is always interpreted
through a physically admissible thickness, even if the head temporarily
overshoots during Newton iterations.

The cell transmissivity is then
\[
T_i(h) = K_i\,b_i(h).
\]

\section{Cell Balance And Sign Convention}

The solver is built around a residual balance written per cell. The convention
used by the code is:

\begin{itemize}
\item positive residual contributions tend to remove water from the cell,
\item recharge and positive well injection add water and therefore appear with a
      minus sign in the residual,
\item a converged solution satisfies $R_i \approx 0$ for every cell $i$.
\end{itemize}

This convention is the one implemented in
\texttt{assemble\_steady\_residual()} and
\texttt{assemble\_transient\_residual()}.

\section{Internal Edge Fluxes}

Consider an internal edge $e$ shared by cells $a$ and $b$. The code defines one
oriented flux $q^{\text{int}}_e$ that is positive when water leaves cell $a$ and
enters cell $b$:
\[
q^{\text{int}}_e = -\tau_e(h_b - h_a).
\]

The edge transmissive factor is
\[
\tau_e
=
K^{H}_e\,\bar{b}_e\,\frac{L_e}{d_e},
\]
where:

\begin{itemize}
\item $K^{H}_e$ is the harmonic mean of $K_a$ and $K_b$,
\item $\bar{b}_e = \frac{1}{2}(b_a + b_b)$ is the averaged saturated thickness,
\item $L_e$ is the edge length,
\item $d_e$ is the centroid-to-centroid distance between the two cells.
\end{itemize}

The residual contribution of internal fluxes is conservative:
\[
R^{\text{int}}_a \mathrel{+}= q^{\text{int}}_e,
\qquad
R^{\text{int}}_b \mathrel{-}= q^{\text{int}}_e.
\]

This part is implemented by:

\begin{itemize}
\item \texttt{internal\_edge\_flux\_from\_head()},
\item \texttt{accumulate\_internal\_flux\_residual()}.
\end{itemize}

\section{Imposed-Head Exchanges}

Some edges carry an imposed stage $H_e$, for example:

\begin{itemize}
\item side Dirichlet conditions,
\item stream supports,
\item ocean supports.
\end{itemize}

For such edges, the code builds cell-to-edge transmissive coefficients
\[
\tau_{a,e} = K_a\,b_a\,\frac{L_e}{d_{a,e}},
\]
and similarly for $\tau_{b,e}$ when a second owner cell exists.

The corresponding exchange fluxes are
\[
q^{D}_{a,e} = -\tau_{a,e}(H_e - h_a),
\qquad
q^{D}_{b,e} = -\tau_{b,e}(H_e - h_b).
\]

Those fluxes are accumulated directly into the owner-cell residuals. In the
current mesh usage, imposed-head support is mainly used on boundary edges, but
the implementation is slightly more general and can also account for an edge
with two owner cells.

\section{Recharge, Wells, Drainage And Saturation Excess}

\subsection{Recharge}

Recharge is represented as a surface rate $r_i$ in m/s. In the current first
slice of the backend, recharge is homogeneous in space for a given stress
period. Its volumetric contribution in cell $i$ is
\[
A_i r_i.
\]

\subsection{Wells}

Wells are converted to cell-based volumetric rates $Q^{\text{well}}_i$ in m$^3$/s.
The sign convention is:

\begin{itemize}
\item positive value: injection into the aquifer,
\item negative value: pumping from the aquifer.
\end{itemize}

Since injection adds water, the residual contains $-Q^{\text{well}}_i$.

\subsection{Drainage}

Drainage is a simple top leakage law that activates only when the head exceeds
the cell top elevation:
\[
q^{\text{drain}}_i = C_i^{\text{drain}}
\max(h_i - z^{\text{top}}_i, 0).
\]

This term is always treated as an outflow from the aquifer and therefore enters
the residual with a positive sign.

\subsection{Regularized Saturation Excess}

The current backend includes a regularized saturation-excess source term. Its
goal is pragmatic: avoid a sharp on/off overflow law exactly at full saturation
while still allowing near-surface water release when the cell is almost full.

Define the saturation ratio
\[
\sigma_i(h) =
\frac{b_i(h)}{z^{\text{top}}_i - z^{\text{bot}}_i}.
\]

Let $R^{\text{lat}}_i$ be the lateral balance rate converted back to a surface
flux by division through $A_i$. The code forms
\[
\beta_i
=
\max\!\left(
    -\frac{R^{\text{lat}}_i}{A_i} + \max(r_i, 0),
    0
\right),
\]
and then applies the regularization
\[
s_i(h)
=
\exp\!\left(
    -\frac{1 - \sigma_i(h)}{\varepsilon}
\right)
\beta_i,
\]
where $\varepsilon > 0$ is the regularization radius.

This term is not meant to be a final physically complete overland-flow model.
It is a smooth release mechanism used by the current backend close to full
saturation.

\section{Steady Residual}

The steady residual assembled by the code is
\[
R_i^{\text{steady}}(h)
=
R_i^{\text{int}}(h)
+
R_i^{D}(h)
+
q_i^{\text{drain}}(h)
+
A_i s_i(h)
-
A_i r_i
-
Q_i^{\text{well}}.
\]

The steady solve therefore searches for
\[
R_i^{\text{steady}}(h) = 0
\qquad \text{for all cells } i.
\]

\section{Transient Residual}

For one fully implicit backward-Euler time step from $t^n$ to $t^{n+1}$, the
additional storage term is
\[
R_i^{\text{storage}}
=
A_i S_i \frac{h_i^{n+1} - h_i^n}{\Delta t}.
\]

The transient residual is then
\[
R_i^{\text{transient}}(h^{n+1})
=
A_i S_i \frac{h_i^{n+1} - h_i^n}{\Delta t}
+
R_i^{\text{int}}(h^{n+1})
+
R_i^{D}(h^{n+1})
+
q_i^{\text{drain}}(h^{n+1})
+
A_i s_i(h^{n+1})
-
A_i r_i^{n+1}
-
Q_{i}^{\text{well},\,n+1}.
\]

The backend solves
\[
R_i^{\text{transient}}(h^{n+1}) = 0
\qquad \text{for all cells } i.
\]

\section{Nonlinear Solution Strategy}

\subsection{Local Runtime}

The local runtime uses a dense damped Newton loop:

\begin{enumerate}
\item start from an initial head guess,
\item assemble the residual,
\item build a dense Jacobian by forward finite differences,
\item solve the linearized Newton system,
\item backtrack the Newton update until the residual decreases,
\item stop when the infinity norm of the residual is below tolerance.
\end{enumerate}

This strategy is transparent and easy to debug, which is why it is useful in an
early implementation slice.

\subsection{SciPy Runtime}

The SciPy runtime keeps the same residual and the same finite-difference
Jacobian, but delegates the nonlinear solve to
\texttt{scipy.optimize.root}. The important point is that the physics does not
change; only the nonlinear driver changes.

\section{Mapping Between Equations And Code}

\begin{center}
\begin{tabular}{ll}
\toprule
Mathematical object & Main implementation point \\
\midrule
$b_i(h)$ & \texttt{saturated\_thickness\_from\_head()} \\
$T_i(h)$ & \texttt{transmissivity\_from\_head()} \\
$q^{\text{int}}_e$ & \texttt{internal\_edge\_flux\_from\_head()} \\
$q^{D}_e$ & \texttt{imposed\_head\_edge\_flux\_from\_head()} \\
$q_i^{\text{drain}}$ & \texttt{drainage\_outflow\_from\_head()} \\
$s_i(h)$ & \texttt{saturation\_excess\_rate\_from\_balance()} \\
$R_i^{\text{steady}}$ & \texttt{assemble\_steady\_residual()} \\
$R_i^{\text{transient}}$ & \texttt{assemble\_transient\_residual()} \\
runtime solve & \texttt{local\_runtime.py} / \texttt{scipy\_runtime.py} \\
problem orchestration & \texttt{boussinesq.py} \\
\bottomrule
\end{tabular}
\end{center}

\section{Interpretation And Current Limits}

The current backend is best viewed as a clear and auditable first finite-volume
Boussinesq implementation, not yet as the final production solver. In
particular:

\begin{itemize}
\item the Jacobian is still dense,
\item the overflow-related term is regularized rather than derived from a full
      coupled surface-flow model,
\item the current slice focuses on small validation meshes and controlled test
      cases.
\end{itemize}

That is acceptable at this stage because the code is optimized for
interpretability, validation and iterative refinement of the physics.

\section{Compilation Hint}

If a LaTeX distribution is installed, this note can usually be compiled from
the package directory with either:

\begin{verbatim}
pdflatex -interaction=nonstopmode boussinesq_math_notes.tex
\end{verbatim}

or:

\begin{verbatim}
latexmk -pdf boussinesq_math_notes.tex
\end{verbatim}

\end{document}
